\section{Tangent boxes}\label{sec:boxes}

Throughout this section, \(X\) is a smooth projective complex manifold
with a fixed splitting \(T_X=E_1\oplus E_2\).
Let \(P_i:T_X\to E_i\subset T_X\) be the corresponding holomorphic
projectors. Put
\[
C=\PP^1,\qquad S=C\times C,\qquad
\Delta=\{(s,s):s\in C\},
\]
and write \(T^{(i)}=\pr_i^*T_C\subset T_S\).

\begin{proposition}[Tangent box]\label[proposition]{prop:box}
Every holomorphic map \(g:C\to X\) admits a holomorphic map \(F:S\to X\)
such that
\begin{equation}\label{eq:box}
F|_\Delta=g,\qquad
dF(T^{(1)})\subset F^*E_1,\qquad
dF(T^{(2)})\subset F^*E_2.
\end{equation}
Neither summand is assumed integrable.
\end{proposition}

The proof has three stages. The projectors first determine a formal
map along \(\Delta\). A count of formal sections then produces a
rational map with that expansion near part of the diagonal. Finally,
the two tangency conditions force every exceptional curve in a
resolution to be contracted, so the rational map is holomorphic on
all of \(S\).

\subsection{The formal Cauchy problem}

We specify the analytic formal neighborhood in which the first stage
takes place. Let \(\mathcal I\subset\OO_S\) be the analytic ideal
sheaf of \(\Delta\), and let \(\iota:\Delta\hookrightarrow S\) be
the inclusion. On the topological space \(\Delta\), set
\[
\wh{\OO}_{S,\Delta}
 =\varprojlim_{n\ge0}\iota^{-1}(\OO_S/\mathcal I^{n+1}),
\qquad
\wh{\mathcal I}=\ker(\wh{\OO}_{S,\Delta}\longrightarrow\OO_\Delta).
\]
Thus the analytic completion is the ringed space
\(\wh S_\Delta=(\Delta,\wh{\OO}_{S,\Delta})\). On a diagonal
coordinate disk with coordinates \(u=s\), \(v=t-s\), its structure
sheaf is \(\OO_\Delta[[v]]\), and \(\wh{\mathcal I}=(v)\).
The coefficients are holomorphic along the disk; convergence in the
normal variable \(v\) is not required.

A formal holomorphic map \(\wh G:\wh S_\Delta\to Y\) to a complex
manifold \(Y\) consists of a holomorphic diagonal map
\(g_0:\Delta\to Y\) and a morphism of sheaves of \(\CC\)-algebras
\[
\wh G^\#:g_0^{-1}\OO_Y\longrightarrow\wh{\OO}_{S,\Delta}
\]
reducing to \(g_0^\#\) modulo \(\wh{\mathcal I}\). We require this
map to be given locally by holomorphic Taylor substitution, as follows.
On a source disk with the above \(\OO_\Delta[[v]]\) description,
view functions of \(u\) as constant series in \(v\).
In target coordinates \(z=(z^1,\ldots,z^q)\), put
\(f=(\wh G^\#z^1,\ldots,\wh G^\#z^q)\) and
\(\eta=f-z\circ g_0\in\wh{\mathcal I}^{\oplus q}\). For a
holomorphic function \(h\) on the target chart, substitution at order
\(n\) is
\[
\wh G^\#h\bmod\wh{\mathcal I}^{n+1}
 =\sum_{|\alpha|\le n}
   \frac{(\partial^\alpha h)(z\circ g_0)}{\alpha!}\,
   \eta^\alpha\bmod\wh{\mathcal I}^{n+1}.
\]
Every monomial in \(\eta\) of total degree \(n+1\) is zero at this
order, so the sum is finite.
These substitutions are compatible as \(n\) varies and preserve
products and holomorphic coordinate changes. Consequently, one can
specify a formal map by coordinate series with the stated constant
terms, provided that on overlaps they agree by these substitutions.
The same rule then determines the pullback of every holomorphic
function, not only of the coordinate functions.

A source holomorphic vector field \(D\) satisfies
\(D(\mathcal I^{k+1})\subset\mathcal I^k\) for \(k\ge0\), and
hence acts as a continuous derivation on the completion. Locally,
\(\partial_u\) differentiates the coefficients of a series and
\(\partial_v\) differentiates its powers of \(v\). A derivation
can lower the normal order, so it need not define a derivation on any one fixed-order
thickening. Formal differential identities below are equalities in
the completion. In target coordinates they are equalities of vectors
of formal series; the formal chain rule makes them independent of the
target coordinates.

\begin{lemma}\label[lemma]{lem:formal}
There is a unique formal holomorphic map \(\wh F\) along \(\Delta\)
with diagonal value \(g\) and the two tangency properties in
\eqref{eq:box}, interpreted formally.
\end{lemma}

\begin{proof}
Use the same coordinate on a source chart in each factor, denoting
the two coordinates by \(s,t\). Set
\[
u=s,\qquad v=t-s.
\]
After restricting the diagonal disk so that \(g\) takes values in
one target coordinate chart, impose
\begin{equation}\label{eq:cauchy}
\wh F(u,0)=g(u),\qquad
\partial_v\wh F=P_2(\wh F)\partial_u\wh F.
\end{equation}
Here the entries of the holomorphic matrix \(P_2\) are substituted
by the Taylor rule above. That rule preserves \(P_2^2=P_2\) and
\(P_1+P_2=\id\).

Because \(\partial_u=\partial_s+\partial_t\) and
\(\partial_v=\partial_t\), \eqref{eq:cauchy} is equivalent to
\[
P_1(\wh F)\partial_t\wh F=0,\qquad
P_2(\wh F)\partial_s\wh F=0.
\]
Indeed, the equation puts \(\partial_v\wh F\) in the image of
\(P_2(\wh F)\), while applying \(P_2(\wh F)\) to
\(\partial_u\wh F-\partial_v\wh F\) gives zero. Conversely,
these two statements recover \eqref{eq:cauchy}. They are exactly
the two formal tangency conditions.

Write \(\wh F=\sum_{k\ge0}f_k(u)v^k\), with \(f_0=g\) in the
chosen coordinates. At order zero, \eqref{eq:cauchy} gives
\[
f_1=P_2(g)g'.
\]
Thus along the diagonal the second-factor derivative is
\(P_2(g)g'\), and the first-factor derivative is its complement
\(g'-f_1=P_1(g)g'\). The higher coefficients extend this forced
first-order decomposition. The coefficient of \(v^k\) on the
right side of \eqref{eq:cauchy} involves only
\(f_0,\ldots,f_k\) and their \(u\)-derivatives. It determines
\((k+1)f_{k+1}\), so induction gives a unique solution whose
coefficients are holomorphic on the diagonal disk.

We verify that these local solutions define a formal map globally.
Under a target-coordinate change \(z'=\psi(z)\), Taylor substitution
gives \(\wh F'=\psi(\wh F)\), and the formal chain rule gives
\(D\wh F'=d\psi(\wh F)D\wh F\). The matrix of \(P_2\) changes
by the corresponding conjugation, so \eqref{eq:cauchy} is preserved.
For a source change \(s'=h(s)\), \(t'=h(t)\), one has
\[
u'=h(u),\qquad v'=h(u+v)-h(u).
\]
The new normal parameter is \(v\) times a unit, so this substitution
identifies the two completions. With
\[
a=\frac1{h'(u)},\qquad b=\frac1{h'(u+v)},
\]
the continuous source derivations transform as
\[
\partial_{v'}=b\partial_v,\qquad
\partial_{u'}=a\partial_u+(b-a)\partial_v.
\]
Since \(P_2(\wh F)\partial_v\wh F=\partial_v\wh F\), these
formulas imply
\[
P_2(\wh F)\partial_{u'}\wh F
 =b\partial_v\wh F=\partial_{v'}\wh F.
\]
The initial value also transforms to the same diagonal map. Local
uniqueness therefore identifies the coordinate series on every
overlap. Their Taylor substitution maps agree as well, since those
maps are determined by the coordinate series. This proves both
global existence and uniqueness.
\end{proof}

\subsection{Rational reconstruction along the diagonal}

We next recover an ordinary rational map from the formal one. The
positive normal bundle of \(\Delta\) bounds the number of formal
sections after twisting; comparison with pairs of polynomial sections
then forces each formal coordinate ratio to be rational. The
classical study of algebraic formal rational functions provides
background for this mechanism; see Hironaka--Matsumura \cite{HM68}
and B\u{a}descu \cite{Bad96}. The proof below works
directly in the analytic completion just defined, including the
passage to an ordinary holomorphic map near the diagonal.

We use the notation
\(\OO_S(m,m)=\pr_1^*\OO_C(m)\otimes\pr_2^*\OO_C(m)\).

\Needspace{7\baselineskip}
\begin{lemma}\label[lemma]{lem:rational}
Let \(Y\subset\PP^N\) be a closed complex submanifold and let
\(\wh G\) be a formal holomorphic map from the formal neighborhood
of \(\Delta\subset S\) to \(Y\).
There is a rational map \(R:S\dashrightarrow\PP^N\) that is
holomorphic on a neighborhood of a nonempty open subset of
\(\Delta\), where its formal expansion agrees with \(\wh G\).
On that neighborhood it takes values in \(Y\).
\end{lemma}

\begin{proof}
Write \(g_0=\wh G|_\Delta\). We first construct the formal pullback
of \(\OO_{\PP^N}(1)\). Choose a cover of \(g_0(\Delta)\) by
affine projective charts \(U_a\) with holomorphic frames \(e_a\)
of this line bundle, and write \(e_b=t_{ab}e_a\) on overlaps.
Restricting \(t_{ab}\) to \(Y\) and applying \(\wh G^\#\) gives
formal functions \(\wh t_{ab}\) on
\(g_0^{-1}(U_a\cap U_b)\). Their constant terms are
\(t_{ab}\circ g_0\), which are nowhere zero. Thus each
\(\wh t_{ab}\) is a unit in the completion. The substitutions
preserve the transition cocycle identities, so the relations
\(\wh e_b=\wh t_{ab}\wh e_a\) glue free rank-one modules into a
locally free \(\wh{\OO}_{S,\Delta}\)-module \(L\).

If the homogeneous coordinate section \(x_\nu\) is
\(c_{\nu,a}e_a\) on \(U_a\), evaluate \(c_{\nu,a}|_Y\) by the
same substitution. The resulting local sections
\(\wh G^\#(c_{\nu,a}|_Y)\wh e_a\) agree on overlaps, giving global
sections \(S_0,\ldots,S_N\) of this very module \(L\).
Its reduction \(L_0=L/\wh{\mathcal I}L\) is the ordinary
holomorphic line bundle \(g_0^*\OO_{\PP^N}(1)\) on \(\Delta\simeq C\).
The reductions of the coordinate sections have no common zero, so
after reordering, \(S_0|_\Delta\) is not identically zero. Its zero
divisor has some degree \(d\ge0\) and identifies
\[
L_0\simeq\OO_C(d).
\]
For example, a degree-\(d\) homogeneous polynomial with the same
zero divisor gives this identification directly.

For a fixed \(m\ge0\), let \(\wh{\OO}_S(m,m)\) denote the restriction of
the ordinary line bundle to the completion, and put
\[
\mathcal M_m=L\otimes_{\wh{\OO}_{S,\Delta}}\wh{\OO}_S(m,m),
\qquad V=H^0(\Delta,\mathcal M_m),
\]
\[
V_k=H^0(\Delta,\wh{\mathcal I}^{\,k}\mathcal M_m)\subset V
\quad(k\ge0).
\]
Thus \(V_0=V\), and \(V_k\) consists precisely of global formal
sections vanishing to normal order at least \(k\). Passing a section
to its class modulo \(\wh{\mathcal I}^{k+1}\mathcal M_m\) has
kernel \(V_{k+1}\). It gives the leading-coefficient injection
\begin{equation}\label{eq:graded}
V_k/V_{k+1}\hookrightarrow
H^0\bigl(C,\OO_C(d+2m-2k)\bigr).
\end{equation}
Indeed, the leading-coefficient sheaf is
\[
\frac{\wh{\mathcal I}^{\,k}\mathcal M_m}
     {\wh{\mathcal I}^{\,k+1}\mathcal M_m}
\simeq L_0\otimes\OO_S(m,m)|_\Delta
             \otimes\iota^{-1}(\mathcal I^k/\mathcal I^{k+1}).
\]
Only the constant terms of the frame transitions act on this
quotient. Since \(\Delta\) has bidegree \((1,1)\), its conormal
bundle \(\mathcal I/\mathcal I^2\) is \(\OO_C(-2)\); the displayed
sheaf is therefore \(\OO_C(d+2m-2k)\), as asserted.

This filtration is separated. In a formal frame a section is a
series in \(v\), so membership in every \(V_k\) makes every
coefficient zero. For all sufficiently large \(k\), the space on
the right side of \eqref{eq:graded} is zero. Hence, for some \(K\),
\(V_k=V_{k+1}\) for every \(k\ge K\), and
\(V_K=\bigcap_{k\ge K}V_k=0\) by separatedness. The finite remaining
filtration gives
\begin{equation}\label{eq:sectionbound}
\dim_\CC V
\le\sum_{k\ge0}\max\{d+2m-2k+1,0\}
=m^2+O(m+1).
\end{equation}
Here \(d\) is fixed. More explicitly, the sum is
\((m+q+1)^2\) for \(d=2q\) and
\((m+q+1)(m+q+2)\) for \(d=2q+1\).

Restrict ordinary polynomial sections of \(\OO_S(m,m)\) to the
completion. For each \(\nu=1,\ldots,N\), consider
\[
H^0(S,\OO_S(m,m))^{\oplus2}\longrightarrow V,\qquad
(A,B)\longmapsto AS_0+BS_\nu.
\]
Its source has dimension \(2(m+1)^2\). Choose \(m\) large enough
that this exceeds the bound in \eqref{eq:sectionbound}, and choose
a nonzero kernel pair \((A_\nu,B_\nu)\) for each \(\nu\).

Fix a connected diagonal coordinate disk \(D\) on which
\(S_0|_\Delta\) is nowhere zero, contained in one product affine
chart. The section \(S_0\) is then a formal frame for \(L\) along
\(D\), and \(S_\nu/S_0\) is the substituted affine coordinate
\(x_\nu/x_0\). Necessarily \(B_\nu\ne0\): otherwise \(A_\nu S_0=0\)
would say that the ordinary section \(A_\nu\) has zero formal
expansion along \(D\). Its convergent local Taylor series would
make it zero on a neighborhood of \(D\), and the identity principle
on connected \(S\) would give \(A_\nu=0\), a contradiction.
The ratios
\[
q_\nu=-A_\nu/B_\nu\in\CC(S)
\]
therefore define a rational map \(R=[1:q_1:\cdots:q_N]\).

It remains to show that these ratios are holomorphic simultaneously
near one point of \(D\). Use \(u=s\), \(v=t-s\) and an ordinary
frame for \(\OO_S(m,m)\) on the chosen product chart, and represent
\(A_\nu,B_\nu\) by holomorphic functions \(a_\nu(u,v)\),
\(b_\nu(u,v)\). Write their normal Taylor coefficients on \(D\)
as \(a_{\nu,k}(u)\), \(b_{\nu,k}(u)\). For each \(\nu\), the
generic normal order
\[
r_\nu=\min\{k\ge0:b_{\nu,k}\text{ is not identically zero on }D\}
\]
is finite. Otherwise Taylor uniqueness near \(D\), followed by
the identity principle on \(S\), would make the nonzero ordinary
section \(B_\nu\) vanish identically. Each leading coefficient
\(b_{\nu,r_\nu}\) is a nonzero holomorphic function on the connected
disk \(D\). Its zeros are discrete, so a single point \(p\in D\)
can be chosen where all the finitely many leading coefficients are
nonzero.

In the formal frame \(S_0\), the kernel relations become
\[
a_\nu+b_\nu(S_\nu/S_0)=0
\quad\text{in }\OO_D[[v]].
\]
By the definition of \(r_\nu\), all coefficients of \(b_\nu\)
below that order vanish on \(D\); the relation makes the same true
of \(a_\nu\). Since \(a_\nu\) and \(b_\nu\) are ordinary
holomorphic functions, this formal vanishing implies actual
holomorphic divisibility near \(p\):
\[
a_\nu=v^{r_\nu}\widetilde a_\nu,\qquad
b_\nu=v^{r_\nu}\widetilde b_\nu,
\qquad
\widetilde b_\nu(p)=b_{\nu,r_\nu}(p)\ne0.
\]
Indeed, their convergent Taylor expansions in \(v\) have the first
\(r_\nu\) coefficients zero. Shrink to one neighborhood of \(p\)
where every \(\widetilde b_\nu\) is a holomorphic unit.
Multiplication by \(v\) is injective on formal series, so the same
power can be canceled in each formal relation. It follows that
\[
q_\nu=-\widetilde a_\nu/\widetilde b_\nu
\]
is holomorphic there and has formal expansion \(S_\nu/S_0\).
Thus \(R\) is holomorphic on this common neighborhood. Its affine
coordinate series agree with those of the composition of \(\wh G\)
with \(Y\hookrightarrow\PP^N\). The Taylor substitution rule then
makes these formal maps agree on every holomorphic function on the
projective target chart.

Finally, \(R(p)=g_0(p)\), so shrink the source neighborhood until
its image lies in a target chart with holomorphic defining equations
for \(Y\).
Their compositions with \(R\) have zero normal Taylor series,
because their substitutions by \(\wh G\) are zero. The convergent
Taylor series of these ordinary holomorphic functions make them
vanish on a neighborhood of \(p\). There \(R\) takes values in
\(Y\). In local submanifold coordinates, every holomorphic function
on \(Y\) extends to the ambient chart. The ambient formal equality
therefore also gives equality of the formal maps to \(Y\), as required.
\end{proof}

\begin{remark}\label[remark]{rem:g3}
The normal-bundle calculation in \eqref{eq:sectionbound} is the
rationality mechanism in the proof. The related algebraic language
uses the \(G_3\) property: the restriction map from rational
functions on the ambient variety to formal rational functions along
the subvariety is an isomorphism. The proof above does not invoke a
\(G_3\) theorem. Its section count and cancellation work directly
with the given holomorphic formal series, without an
analytic-to-algebraic comparison.
\end{remark}

\subsection{Complementary tangency removes indeterminacy}

Rational reconstruction has produced a map only away from its base
points. We next use the specified splitting to show that no base point
can remain: each exceptional curve would have to map tangentially to
both complementary summands, forcing its image to be a point.

\begin{lemma}\label[lemma]{lem:removal}
Let \(R:S\dashrightarrow X\) be a rational map such that, on its
domain of definition,
\[
dR(T^{(1)})\subset R^*E_1,\qquad
dR(T^{(2)})\subset R^*E_2.
\]
Then \(R\) extends to a holomorphic map \(S\to X\).
\end{lemma}

\begin{proof}
Choose a projective embedding \(X\subset\PP^N\).
By \cref{lem:surface}, there are a finite set \(B\subset S\), a
composition of point blowups \(\mu:S'\to S\) over \(B\), and a
holomorphic map \(F':S'\to\PP^N\) extending \(R\circ\mu\).
Since \(X\) is closed and \(F'\) maps a dense open set into \(X\),
the entire image lies in \(X\).

We show that \(F'\) is constant on each exceptional component.
Let \(H\) be such a component, lying over \(p\in B\), and let
\(s,t\) be coordinates on the two source factors near \(p\).
On the inverse image of this coordinate neighborhood,
\(P_i dF'\) is a holomorphic one-form with values in \(F'^*T_X\).
The tangency assumptions imply the identities
\begin{equation}\label{eq:wedges}
(P_1dF')\wedge d(s\circ\mu)=0,\qquad
(P_2dF')\wedge d(t\circ\mu)=0.
\end{equation}
The two wedge expressions are sections of the locally free sheaf
\(\Omega^2_{S'}\otimes F'^*T_X\). The identities first hold off
the exceptional curves, where \(\mu\) is an isomorphism, and then
hold everywhere by holomorphicity.

Both coordinate functions \(s\circ\mu-s(p)\) and
\(t\circ\mu-t(p)\) vanish along \(H\), since \(\mu(H)=\{p\}\).
Their generic orders along \(H\) are positive and finite. Indeed,
in a normal coordinate along a connected disk of \(H\), the first
normal Taylor coefficient that is not identically zero gives that
order. If no such coefficient existed, the pullback would vanish on
an open neighborhood meeting the dense locus where \(\mu\) is an
isomorphism, which is impossible for a coordinate function. Denote
the two respective orders by \(\alpha,\beta\). This reasoning
applies also to components created by later blowups.
At a general point of \(H\), avoiding intersections with other
components and zeros of the two leading coefficients, choose a
local equation \(r=0\). We can then write
\[
s\circ\mu-s(p)=r^\alpha h,\qquad
t\circ\mu-t(p)=r^\beta\widetilde h
\]
with holomorphic units \(h,\widetilde h\).

The first identity in \eqref{eq:wedges} is therefore
\[
r^{\alpha-1}(P_1dF')\wedge(\alpha h\,dr+r\,dh)=0.
\]
The holomorphic local ring of the smooth surface is a domain, so
multiplication by \(r^{\alpha-1}\) is injective on the locally free
module of bundle-valued two-forms. It can therefore be canceled.
At \(r=0\),
evaluate the resulting ambient two-form on a vector \(w\) tangent
to \(H\) and a vector \(n\) with \(dr(n)=1\). The result is
\(\alpha h\,P_1dF'(w)=0\), hence \(P_1dF'(w)=0\).
Canceling \(r^{\beta-1}\) in the second identity gives
\(P_2dF'(w)=0\) in the same way. Since \(P_1+P_2=\id\), the
derivative of \(F'|_H\) vanishes on a dense open subset of \(H\),
and hence everywhere. Since \(H\) is connected, \(F'|_H\) is constant.

The fibers of \(\mu\) away from \(B\) are points. Each nontrivial
fiber over \(B\) is a connected union of exceptional components:
the first blowup creates an exceptional line, and every later blowup
over that point replaces its center by a new line meeting the strict
transforms of the components through it, as in \cref{lem:surface}.
The constants on intersecting components agree,
so \(F'\) is constant on every fiber.

The map \(\mu\) is a continuous surjection from a compact space to
a Hausdorff space. It is therefore a quotient map, and \(F'\)
descends to a continuous map \(F:S\to X\), holomorphic off \(B\).
Near each point of \(B\), continuity places the image in a single
target coordinate chart. The coordinate functions are continuous
and holomorphic away from the isolated point. One-variable
removability on coordinate lines makes them separately
holomorphic, and the iterated Cauchy formula makes them
holomorphic. Hence \(F\) is holomorphic everywhere.
\end{proof}

\begin{proof}[Proof of \cref{prop:box}]
Apply \cref{lem:formal} and then \cref{lem:rational} to a projective
embedding of \(X\). The resulting rational map is holomorphic near
a nonempty diagonal patch and has formal expansion \(\wh F\) there.
Continuous differentiation and Taylor substitution identify the
expansions of \(P_1(R)\partial_tR\) and \(P_2(R)\partial_sR\)
with the zero formal tangency expressions of \(\wh F\). These are
ordinary holomorphic expressions, so Taylor uniqueness makes them
zero on a neighborhood of that patch. Let \(B\) be the finite base set from
\cref{lem:surface}. The complement \(S\setminus B\) is connected;
one can perturb any path in the connected real four-manifold \(S\)
inside coordinate balls to avoid the finitely many points.
The inverse image of \(X\) under the holomorphic map
\(R:S\setminus B\to\PP^N\) is an analytic subset containing a
nonempty open set, so it is all of \(S\setminus B\). The two
tangency conditions are then vanishing identities for holomorphic
bundle morphisms on this connected complement, and continue there
by the identity principle.

Now \cref{lem:removal} gives a holomorphic extension \(F:S\to X\).
Tangency extends by holomorphicity. The maps \(F|_\Delta\) and
\(g\) agree on the original nonempty diagonal patch, so they
agree on all of \(\Delta\).
\end{proof}
